\documentclass{article} % This command is used to set the type of
% document you are working on such as an article, book, or presenation

\usepackage{
  amsmath,      % Math Environments
  amssymb,      % Extended Symbols
  enumerate,        % Enumerate Environments
  graphicx,      % Include Images
  lastpage,      % Reference Lastpage
  multicol,      % Use Multi-columns
  multirow,      % Use Multi-rows
  cases,
}

\usepackage{geometry} % This package allows the editing of the page layout
\usepackage{amsmath}  % This package allows the use of a large range
% of mathematical formula, commands, and symbols
\usepackage{graphicx}  % This package allows the importing of images

\newcommand{\question}[2][]{
  \begin{flushleft}
    \textbf{Question #1}: \textit{#2}

\end{flushleft}}
\newcommand{\sol}{\textbf{Solution}:} %Use if you want a boldface solution line
\newcommand{\maketitletwo}[2][]{
  \begin{center}
    \Large{\textbf{Assignment #1}

    Mathematical Statistics} % Name of course here
    \vspace{5pt}

    \normalsize{N. Ohayon Rozanes  % Your name here

    \today}        % Change to due date if preferred
    \vspace{15pt}

\end{center}}
\begin{document}
\maketitletwo[2]  % Optional argument is assignment number
%Keep a blank space between maketitletwo and \question[1]

\question[1]{ }

\begin{enumerate}[(a)]
  \item $Y = X_1 + X_2$ Since $X_1$ and $X_2$ are independent, we can
    use the formula
    \[
      P(Y = y) = \sum_{k=-\infty}^\infty P(X_1 = k)P(X_2 = y - k)
    \]
    Since the support of a poisson distribution is $k \ge 0$ this
    forces $y \ge k$, thus the sum becomes and evaluates to
    \begin{align*}
      P(Y = y) &= \sum_{k=0}^y P(X_1 = k)P(X_2 = y - k)\\
      &= \sum_{k=0}^y
      \frac{\lambda^k\lambda_2^{y-k}}{k!(y-k)!}e^{-(\lambda_1 + \lambda_2)}\\
      &= \sum_{k=0}^y \frac{1}{y!}\binom{y}{k}\lambda^k_1
      \lambda_2^{y-k}e^{-(\lambda_1 + \lambda_2)}\\
      &= \frac{1}{y!}e^{-(\lambda_1 + \lambda_2)}
      \sum_{k=0}^y
      \binom{y}{k}\lambda^k_1
      \lambda_2^{y-k}\\
    \end{align*}
    By way of the binomial theorem
    \[
      \sum_{k=0}^y
      \binom{y}{k}\lambda^k_1
      \lambda_2^{y-k} =
      (\lambda_1 + \lambda_2)^y
    \]
    Thus
    \[
      P(Y = y) = \frac{(\lambda_1 + \lambda_2)^y}{y!}e^{-(\lambda_1 +
      \lambda_2)}
    \]
  \item The above PMF follows exactly the Poisson$(\lambda_1 +
    \lambda_2)$ distribution.
  \item The mgf of $Y = X_1 + X_2$ is
    $$
    \text{mgf}_Y(t) = \mathbb{E}(e^{tY}) = \mathbb{E}(e^{tX_1 + tX_2})
    $$

    Since $X_1$ and $X_2$ are independent we can say that
    $$
    \text{mgf}_{X_1}(t)\text{mgf}_{X_2}(s) = \text{mgf}_{X_1, X_2}(t,
    s) = \mathbb{E}(e^{tX_1 + sX_2})
    $$
    Thus
    \[
      \text{mgf}_{Y}(t) = \text{mgf}_{X_1, X_2}(t, t)
    \]
    However:
    \[
      \text{mgf}_{X_1, X_2}(t,t) =
      \text{mgf}_{X_1}(t)\text{mgf}_{X_2}(t)=
      \mathbb{E}(e^{tX_1})\mathbb{E}(e^{tX_2})
    \]
    The right hand side calculates to
    $$
    \mathbb{E}(e^{tX_1})\mathbb{E}(e^{tX_2}) =
    e^{\lambda_1(e^t-1)}e^{\lambda_2(e^t-1)} = e^{(e^t-1)(\lambda_1 +\lambda_2)}
    $$
    Thus,
    \[
      \text{mgf}_Y(t) =  e^{(e^t-1)(\lambda_1 +\lambda_2)}
    \]
    Which is exactly the MGF of Poisson$(\lambda_1 +\lambda_2)$, thus
    Y follows a Poisson$(\lambda_1 + \lambda_2)$ distribution.

\end{enumerate}

\question[2]{ }
\begin{enumerate}[(a)]
  \item The support of $X$, $\mathbb{X}$ is simply the set $\{(x_1,
    x_2): 0 < x_1^2 + x_2^2 < 1 \text{ and } x_1\neq0, x_2\neq0\}$ as
    given by the indicator function. Which is The unit disk, without
    the unit axis and origin.

    Similarily, the support of $Y_1$ is given almost immediatly by
    the support of $X$, that is $0 < x_1^2 + x_2^2 < 1 \implies 0 <
    y_1 < 1$ thus the support of $Y_1$ is $(0,1)$. Given that, it is
    clear to see that the range of $Y_2$ is also $(0,1)$. Meaning
    that the support of $Y$, $\mathbb{Y}$ is the unit square in
    $\mathbb{R}^2$ without the borders.
  \item
    \begin{align*}
      Y_1 = X_1^2 + X_2^2 \implies Y_2 &= \frac{X_1^2}{Y_1} \\
      Y_2Y_1 &= X_1^2 \\
      \pm\sqrt{Y_2Y_1} &= X_1
    \end{align*}
    Given that
    \begin{align*}
      Y_1 &= (\pm\sqrt{Y_2Y_1})^2 + X_2^2 \\
      Y_1 &= Y_2Y_1 + X_2^2\\
      \pm\sqrt{Y_1(1 - Y_2)} = X_2
    \end{align*}

    Yielding the answer
    \[
      \begin{pmatrix}
        X_1 \\ X_2
      \end{pmatrix}
      =
      \begin{pmatrix}
        \pm\sqrt{Y_2Y_1} \\
        \pm\sqrt{Y_1(1-Y_2)}
      \end{pmatrix}
    \]

    Since for any $(Y_1, Y_2)$, there are four possible $(X_1,
    X_2)$ vectors which might have produced it, $u$ is a $4$-to-$1$ function.
  \item Since the magnitude is the same between signs, we can choose
    one branch and the answer will be the same
    \[
      |\text{Det}(u)| =
      \begin{vmatrix}
        \dfrac{y_2}{2\sqrt{y_1y_2}} &
        \dfrac{y_1}{2\sqrt{y_1y_2}}
        \\[10pt]
        \dfrac{1-y_2}{2\sqrt{y_1(1-y_2)}} &
        -\dfrac{y_1}{2\sqrt{y_1(1-y_2)}}
      \end{vmatrix}
    \]
    Which simplifies down to:
    \[
      \frac{-y_1}{4\sqrt{y_1y_2}\sqrt{y_1(1-y_2)}} =
      \frac{1}{4\sqrt{y_2(1-y_2)}}
    \]
  \item
    Using this,we can calculate the joint pdf of $Y$
    \[
      f_Y(y_1, y_2) = f_X(\pm\sqrt{y_2y_1},
      \pm\sqrt{y_1(1-y_2)})\frac{1}{4\sqrt{y_2(1-y_2)}} \cdot 4
    \]
    Where the $4$ coefficient comes from the $4$ preimages, Then
    integrating away the $y_2$ term:
    \begin{align*}
      \int_0^1 \frac{1}{\pi} \mathbf{1}_{\{0 < y_2y_1 + y_1(1-y_2) <
      1\}}\frac{1}{\sqrt{y_2(1-y_2)}} dy_2
    \end{align*}
    notice that due to the squaring, we can again ignore the $\pm$
    branches. We can simplify the range inside the indicator function
    to just $0 < y_1 < 1$. Thus the indicator function is independent
    of $y_2$, so we can pull it out of the integral
    \begin{align*}
      \frac{1}{\pi} \mathbf{1}_{\{0 < y_1 <1\}}\int_0^1
      \frac{1}{\sqrt{y_2(1-y_2)}} dy_2 &=
      \frac{1}{\pi} \mathbf{1}_{\{0 < y_1 <1\}}\int_0^1
      \frac{1}{}y_2^{-\frac12}(1-y_2)^{-\frac12} dy_2
    \end{align*}
    Is the beta integral $B(\frac{1}2, \frac{1}2) = \pi$. Thus:
    \[
      f_{Y_1}(y_1) = \mathbf{1}_{\{0 < y_1 < 1\}}
    \]
    Which is the uniform distribution from $(0,1)$. A similar for
    integration for $Y_2$ follows
    \begin{align*}
      \int_0^1 \frac{1}{\pi} \mathbf{1}_{\{0 < y_1 <
      1\}}\frac{1}{\sqrt{y_2(1-y_2)}} dy_1 &=
      \frac{1}{\pi\sqrt{y_2(1-y_2)}}\int_0^1  \mathbf{1}_{\{0 < y_1 <
      1\}} dy_1
      \\
      &= \frac{1}{\pi}(y_2)^{-\frac{1}2}(1-y_2)^{-\frac{1}{2}}
    \end{align*}
    Which is exactly the Beta$(\frac{1}2, \frac{1}2)$.
\end{enumerate}

\question[3]{ }
For
\[
  Y = u(X) =
  \begin{pmatrix}
    X_1/X_2 // X_2
  \end{pmatrix}
\]
The inverse function is
$$
Y_2 = X_2 \implies Y_1 = X_1/Y_2 \implies Y_1Y_2 = X_1
$$
The Jacobian of the inverse of $u$ is
$$
J(u) =
\begin{pmatrix}
  y_2 & y_1 \\
  0 & 1
\end{pmatrix}
$$
\[
  \left|J(u^{-1})\right| = |y_2|
\]

Therefore the pdf is
\[
  f_Y(y_1, y_2) = f_X(y_1y_2, y_2)y_2 = 8y_2^32y_1\mathbf{1}_{(0 <
  y_1y_2 < y_2 < 1)}
\]
The ranges of the indicator function can be simplified a little bit
by dividing through by $y_2$ which will yield the seperate ranges $0
< y_1 < 1$ and $0 < y_2 < 1$, so to integrate away $y_1$ and $y_2$ to
obtain the respective marginals
\[
  f_{Y_1}(y_1)
  = \int_0^1 f_{Y_1,Y_2}(y_1,y_2)\,dy_2
  = \int_0^1 8y_1 y_2^3 \, dy_2
  = 8y_1 \int_0^1 y_2^3 \, dy_2
  = 8y_1 \left[\frac{y_2^4}{4}\right]_0^1
  = 2y_1, \quad 0<y_1<1.
\]
And
\[
  f_{Y_2}(y_2)
  = \int_0^1 f_{Y_1,Y_2}(y_1,y_2)\,dy_1
  = \int_0^1 8y_1 y_2^3 \, dy_1
  = 8y_2^3 \int_0^1 y_1 \, dy_1
  = 8y_2^3 \left[\frac{y_1^2}{2}\right]_0^1
  = 4y_2^3, \quad 0<y_2<1.
\]

\[
  f_{Y_1}(y_1)f_{Y_2}(y_2)
  = (2y_1)(4y_2^3)
  = 8y_1y_2^3.
\]

Thus the joint pdf is equal to product of the marginals, and we have
verified that the variables are indpednent.

\question[4]{ }
\begin{enumerate}
  \item
    Given that $f_X(x)$ is location scale and symmetric about $0$, then by
    problem 2.26, $f$ has median and expected value $0$. Define
    \[
      Y = \mu + \sigma X
    \]

    Since $X$ is location-scale, $Y$ has the pdf
    \[
      f_Y = \frac{1}{\sigma}f_X\left(\frac{x-\mu}{\sigma}\right)
    \]
    Notice that
    \begin{align*}
      E(Y) &= E(\mu + \sigma X) \\
      &= \mu + \sigma E(X) \\
      &= \mu + 0 \\
      &= \mu
    \end{align*}
    Thus, again by 2.26,
    $\frac{1}{\sigma}f_X\left(\frac{x-\mu}{\sigma}\right)$ has a median at $0$
  \item Compute $P(X\le x_\alpha)$
    \[
      P(X \ge x_\alpha)
      = \int_{x_\alpha}^{\infty}
      \frac{1}{\sigma}
      f\!\left(\frac{x-\mu}{\sigma}\right)
      \,dx.
    \]
    let $z = \frac{x-\mu}{\sigma} \implies dz = \frac{1}{\sigma}$ thus:
    \[
      P(X \ge x_\alpha)
      = \int_{z_\alpha}^{\infty}
      \frac{1}{\sigma} f(z)\,\sigma\,dz
      = \int_{z_\alpha}^{\infty} f(z)\,dz
      = \alpha.
    \]
    As desired.
\end{enumerate}
\question[5] { }
\begin{enumerate}
  \item
    \[
      P(X \le \mu)
      = \int_{-\infty}^{\mu}
      \frac{1}{\sigma\pi\left(1+\left(\frac{x-\mu}{\sigma}\right)^2\right)}\,dx.
    \]
    \[
      \text{Let } z=\frac{x-\mu}{\sigma},\quad dx=\sigma\,dz.
      \text{ When } x=\mu,\ z=0;\ \text{as } x\to -\infty,\ z\to -\infty.
    \]

    \[
      P(X \le \mu)
      = \int_{-\infty}^{0} \frac{1}{\sigma\pi(1+z^2)}\,\sigma\,dz
      = \int_{-\infty}^{0} \frac{1}{\pi(1+z^2)}\,dz.
    \]

    \[
      \int_{-\infty}^{0} \frac{1}{\pi(1+z^2)}\,dz
      = \frac{1}{\pi}\Big[\arctan(z)\Big]_{-\infty}^{0}
      = \frac{1}{\pi}\left(0-\left(-\frac{\pi}{2}\right)\right)
      = \frac{1}{2}.
    \]

    \[
      \Rightarrow\quad P(X\le \mu)=\frac12
      \quad\text{and hence}\quad
      P(X\ge \mu)=1-P(X<\mu)=\frac12.
    \]
  \item We first calculate the case $\mu = 0$ and $\sigma = 1$
    \[
      P(Z \ge 1)
      = \int_{1}^{\infty} \frac{1}{\pi(1+z^2)}\,dz
      = \frac{1}{\pi}\Big[\arctan z\Big]_{1}^{\infty}
      = \frac{1}{\pi}\left(\frac{\pi}{2}-\frac{\pi}{4}\right)
      = \frac{1}{4}.
    \]

    \[
      P(Z \le -1)
      = \int_{-\infty}^{-1} \frac{1}{\pi(1+z^2)}\,dz
      = \frac{1}{\pi}\Big[\arctan z\Big]_{-\infty}^{-1}
      = \frac{1}{\pi}\left(-\frac{\pi}{4}-\left(-\frac{\pi}{2}\right)\right)
      = \frac{1}{4}.
    \]

    % Thus, for the standard Cauchy, 1 and -1 are the upper and lower quartiles:
    \[
      P(Z \ge 1)=\frac14,
      \qquad
      P(Z \le -1)=\frac14.
    \]
    From 3.38, we define $Z = (X -\mu)/\sigma$, therefore for any $z_0$
    $$
    P(X \ge \mu+\sigma z_0) = P(Z \ge z_0)
    $$
    And also
    $$
    P(X \le \mu+\sigma z_0) = P(Z \le z_0)
    $$
    Therefore:
    \[
      P(X \ge \mu+\sigma)
      = P\!\left(\frac{X-\mu}{\sigma} \ge 1\right)
      = P(Z \ge 1)
      = \frac{1}{4}.
    \]

    \[
      P(X \le \mu-\sigma)
      = P\!\left(\frac{X-\mu}{\sigma} \le -1\right)
      = P(Z \le -1)
      = \frac{1}{4}.
    \]

    % Conclusion:
    \[
      \,P(X \ge \mu+\sigma)=\frac14,\qquad P(X \le \mu-\sigma)=\frac14\,
    \]
\end{enumerate}
\end{document}
